\documentclass[12pt]{article}
\usepackage[utf8]{inputenc}
\usepackage[T1]{fontenc}
\usepackage{hyperref}
\title{Moldova Sign Language Language Resources for K-12 Curriculum Development}
\author{Compiled by Claude Sonnet 5 (claude-sonnet-5) for Andr\'as Kornai}
\date{}
\begin{document}
\maketitle

\section{Introduction}
Moldova Sign Language (ISO 639-3: vsi) is the sign language of the deaf community of Moldova. This search found no textbooks, dictionaries or curriculum materials in the language itself. What it found is advocacy reports, news items and one research article about deaf education. They say that sign language is not adequately recognised, that teachers in special schools receive no sign-language training, and that few interpreters are available. Figures for interpreters differ between sources.

\section{Primary Educational Resources}
No primers or textbooks in Moldova Sign Language were found. The following sources describe the school context.
\begin{itemize}
\item News report that only one special school for deaf children, No.\ 12 in Chi\c{s}in\u{a}u, remains, and that integration into mainstream schools has been criticised by the Association of the Deaf: \url{https://www.old.ipn.md/en/integration-of-deaf-children-into-general-education-schools-failed-association-7967_1081845.html}
\item News report on difficulties of hearing-impaired students in the learning process: \url{https://www.old.ipn.md/en/hearing-impaired-students-experience-difficulties-in-learning-process-associatio-7967_1027168.html}
\item Comments from the Association of the Deaf on the draft ``Education 2030'' implementation programme. They propose recognising Moldovan Sign Language as the main language of instruction and teaching it as a subject: \url{https://particip.gov.md/comment_attachments/4395__2599.pdf}
\end{itemize}

\section{Linguistic and Reference Materials}
No dictionary or grammar was found in this search.
\begin{itemize}
\item Research article by Maximciuc and Bargan on sign language in early education: \url{https://asociatia-alpha.ro/Jrls/041-2025/Jrlsd-041-063.pdf}
\item NGO report to the UN Committee on the Rights of the Child on the lack of official recognition of sign language and on teacher training: \url{https://cnpdc.gov.md/sites/default/files/document/attachments/int_crc_ngo_mda_28936_e.pdf}
\end{itemize}

\section{Religious and Community Resources}
\begin{itemize}
\item Joshua Project profile of the Deaf in Moldova (demographic and missions information only): \url{https://www.joshuaproject.net/index.php/people_groups/19007/MD}
\item Association of the Deaf of Moldova: it appears in the news items above, but no website was found in this search.
\end{itemize}

\section{Digital and Audio Resources}
\begin{itemize}
\item Article ``Not many people in Moldova know sign language'': \url{https://ipn.md/en/not-many-people-in-moldova-know-sign-language-7978_1045940.html}
\item Blog on learning sign language and inclusion: \url{https://euneighbourseast.eu/young-european-ambassadors/blog/blog-learning-sign-language-a-huge-step-to-inclusion/}
\end{itemize}
No video corpus or sign-language dictionary was located.

\section{Curriculum Development Context}
Very little material exists in or for the language. A curriculum would have to start from basic documentation: a sign lexicon, teacher training in the language, and bilingual (Moldova Sign Language and Romanian) instruction. The sources above show that the Association of the Deaf is already asking for these.
\end{document}